% GENERATED FILE. Edit canonical lessons, then run npm run generate:books.
% canonical-source-hash: fnv1a64:e8f113a12aa5d6a1
% canonical-lessons: TA-C288-eppatip-poka-ventum, TA-C288-nerakap-ponkal, TA-R288-asking-the-way

\chapter{Asking the Way}
\label{ch:ta-asking-the-way}
\hlchaptermodality{288}

\begin{chapteropening}
\textbf{By the end of this chapter:} \emph{I can ask the way with … eppaṭip pōka vēṇṭum? and eṅkē irukkiṟatu?, ask how far it is, and give directions with nērākap pōṅkaḷ, tirumpuṅkaḷ and landmarks such as kōyilukku etirē.}

Complete the last lesson of chapter 288: ask the way and how far it is, and give the way back with polite requests and landmarks.
\end{chapteropening}

\section[\texorpdfstring{\ta{எப்படிப்} \ta{போக} \ta{வேண்டும்} (eppaṭip pōka vēṇṭum)}{எப்படிப் போக வேண்டும் (eppaṭip pōka vēṇṭum)}]{\ta{எப்படிப்} \ta{போக} \ta{வேண்டும்} (eppaṭip pōka vēṇṭum) — asking the way}

\label{lesson:TA-C288-eppatip-poka-ventum}



\begin{quote}
Say \textquotedblleft{}I must go.\textquotedblright{} (\emph{nāṉ pōka vēṇṭum}.) Then \textquotedblleft{}the left side\textquotedblright{} and \textquotedblleft{}the right side\textquotedblright{}. (\emph{iṭatu}, \emph{valatu}.) Then say \textquotedblleft{}dampness\textquotedblright{}. (\emph{īram}.)
\end{quote}

\subsection*{You'll want to know: asking the way}
\begin{quote}
\textbf{\ta{மன்னிக்கவும்}, \ta{பேருந்து} \ta{நிலையத்துக்கு} \ta{எப்படிப்} \ta{போக} \ta{வேண்டும்}?}
\end{quote}

\begin{quote}
\emph{maṉṉikkavum, pēruntu nilaiyattukku eppaṭip pōka vēṇṭum?} — \textquotedblleft{}Excuse me, how do I get to the bus station?\textquotedblright{}
\end{quote}

Word for word: \textquotedblleft{}to the bus station, how must (one) go\textquotedblright{}. The place takes \textbf{-\ta{க்கு}}: \textbf{\ta{நிலையத்துக்கு}}, \textquotedblleft{}to the station\textquotedblright{}. After \textbf{\ta{எப்படி}} the \textbf{\ta{ப}} of \textbf{\ta{போக}} doubles in writing. When the place itself is what you want, ask where it is: \textbf{\ta{மருத்துவமனை} \ta{எங்கே} \ta{இருக்கிறது}?} (\emph{maruttuvamaṉai eṅkē irukkiṟatu?}, \textquotedblleft{}where is the hospital?\textquotedblright{})

\subsection*{You'll want to know: how far}
\begin{quote}
\textbf{\ta{இங்கிருந்து} \ta{எவ்வளவு} \ta{தூரம்}?}
\end{quote}

\begin{quote}
\emph{iṅkiruntu evvaḷavu tūram?} — \textquotedblleft{}How far is it from here?\textquotedblright{}
\end{quote}

\textbf{\ta{இங்கிருந்து}} (\emph{iṅkiruntu}, \textquotedblleft{}from here\textquotedblright{}): \textbf{\ta{இங்கே}} with \textbf{-\ta{இருந்து}}, \textquotedblleft{}from\textquotedblright{}. The answer: \textbf{\ta{அது} \ta{அருகில்} \ta{இருக்கிறது}} (\emph{atu arukil irukkiṟatu}, \textquotedblleft{}it is near\textquotedblright{}), or \textbf{\ta{அது} \ta{தூரத்தில்} \ta{இருக்கிறது}} (\emph{atu tūrattil irukkiṟatu}, \textquotedblleft{}it is far away\textquotedblright{}). To ask \textbf{\ta{நடந்து} \ta{போகலாமா}?} (\emph{naṭantu pōkalāmā?}, \textquotedblleft{}can I walk there?\textquotedblright{}), use the \textbf{-\ta{லாமா}} of \textbf{\ta{வரலாமா}}.

\subsection*{Your turn}
Say these aloud:

\begin{itemize}
\raggedright
  \item \textquotedblleft{}Excuse me, how do I get to the bus station?\textquotedblright{} — \emph{maṉṉikkavum, pēruntu nilaiyattukku eppaṭip pōka vēṇṭum?}
  \item \textquotedblleft{}Where is the hospital? How far is it from here?\textquotedblright{} — \emph{maruttuvamaṉai eṅkē irukkiṟatu? iṅkiruntu evvaḷavu tūram?}
\end{itemize}

\begin{tcolorbox}[breakable,colback=teal!4,colframe=teal!35!black,title={Before you move on}]
\textquotedblleft{}How do I get to the station?\textquotedblright{} (\textbf{\ta{நிலையத்துக்கு} \ta{எப்படிப்} \ta{போக} \ta{வேண்டும்}?}) \textquotedblleft{}Where is …?\textquotedblright{} (\textbf{… \ta{எங்கே} \ta{இருக்கிறது}?}) \textquotedblleft{}How far?\textquotedblright{} (\textbf{\ta{எவ்வளவு} \ta{தூரம்}?}) \textquotedblleft{}Can I walk?\textquotedblright{} (\textbf{\ta{நடந்து} \ta{போகலாமா}?})
\end{tcolorbox}
\section[\texorpdfstring{\ta{நேராகப்} \ta{போங்கள்} (nērākap pōṅkaḷ)}{நேராகப் போங்கள் (nērākap pōṅkaḷ)}]{\ta{நேராகப்} \ta{போங்கள்} (nērākap pōṅkaḷ) — go straight, turn, it is opposite}

\label{lesson:TA-C288-nerakap-ponkal}



\begin{quote}
Ask the way to the bus station. (\emph{pēruntu nilaiyattukku eppaṭip pōka vēṇṭum?}) Then say \textquotedblleft{}straight ahead\textquotedblright{} and \textquotedblleft{}to turn\textquotedblright{}. (\emph{nērāka}, \emph{tirumpu}.)
\end{quote}

\subsection*{You'll want to know: the polite request, -\ta{உங்கள்}}
\begin{quote}
\textbf{\ta{நேராகப்} \ta{போங்கள்}. \ta{பிறகு} \ta{இடது} \ta{பக்கம்} \ta{திரும்புங்கள்}.}
\end{quote}

\begin{quote}
\emph{nērākap pōṅkaḷ. piṟaku iṭatu pakkam tirumpuṅkaḷ.} — \textquotedblleft{}Go straight on. Then turn left.\textquotedblright{}
\end{quote}

You have said this ending since the early chapters: \textbf{\ta{சொல்லுங்கள்}}, \textquotedblleft{}tell (me)\textquotedblright{}. It is the verb with \textbf{-\ta{உங்கள்}}, the polite way to ask someone to do something: \textbf{\ta{போ}} $\to$ \textbf{\ta{போங்கள்}}, \textbf{\ta{திரும்பு}} $\to$ \textbf{\ta{திரும்புங்கள்}}, \textbf{\ta{வா}} $\to$ \textbf{\ta{வாருங்கள்}}, \textbf{\ta{பார்}} $\to$ \textbf{\ta{பாருங்கள்}}. \textbf{\ta{பக்கம்}} (\emph{pakkam}, \textquotedblleft{}side\textquotedblright{}): \textbf{\ta{இடது} \ta{பக்கம்}}, \textquotedblleft{}(to) the left side\textquotedblright{}.

Do not mix it up with the \textbf{-\ta{உ}} / \textbf{-\ta{இ}} form of the story lessons. \textbf{\ta{திரும்புங்கள்}} asks someone to turn; \textbf{\ta{திரும்பி}} is \textquotedblleft{}having turned\textquotedblright{}, and waits for another verb: \textbf{\ta{இடது} \ta{பக்கம்} \ta{திரும்பி} \ta{நேராகப்} \ta{போங்கள்}} (\emph{iṭatu pakkam tirumpi nērākap pōṅkaḷ}, \textquotedblleft{}turn left and go straight on\textquotedblright{}).

\subsection*{You'll want to know: landmarks}
\begin{quote}
\textbf{\ta{கோயிலுக்கு} \ta{எதிரே} \ta{ஒரு} \ta{வங்கி} \ta{இருக்கிறது}.}
\end{quote}

\begin{quote}
\emph{kōyilukku etirē oru vaṅki irukkiṟatu.} — \textquotedblleft{}Opposite the temple there is a bank.\textquotedblright{}
\end{quote}

The landmark takes \textbf{-\ta{க்கு}}, then the place word: \textbf{\ta{கோயிலுக்கு} \ta{எதிரே}}, \textquotedblleft{}opposite the temple\textquotedblright{}; \textbf{\ta{வங்கிக்கு} \ta{அருகில்}}, \textquotedblleft{}near the bank\textquotedblright{}. For \textquotedblleft{}between\textquotedblright{}, both landmarks take \textbf{-\ta{க்கும்}}, the \textquotedblleft{}and\textquotedblright{} you know: \textbf{\ta{கடைக்கும்} \ta{பள்ளிக்கும்} \ta{இடையில்}} (\emph{kaṭaikkum paḷḷikkum iṭaiyil}, \textquotedblleft{}between the shop and the school\textquotedblright{}).

\begin{quote}
\textbf{\ta{இரண்டாவது} \ta{தெருவில்} \ta{வலது} \ta{பக்கம்} \ta{திரும்புங்கள்}.}
\end{quote}

\begin{quote}
\emph{iraṇṭāvatu teruvil valatu pakkam tirumpuṅkaḷ.} — \textquotedblleft{}Turn right at the second street.\textquotedblright{}
\end{quote}

\textbf{\ta{இரண்டாவது}} (\emph{iraṇṭāvatu}, \textquotedblleft{}second\textquotedblright{}) is \textbf{\ta{இரண்டு}} with \textbf{-\ta{ஆவது}}; \textbf{\ta{தெருவில்}} (\emph{teruvil}) is \textquotedblleft{}at the street\textquotedblright{}.

\subsection*{Your turn}
Say these aloud:

\begin{itemize}
\raggedright
  \item \textquotedblleft{}Go straight on. Then turn left.\textquotedblright{} — \emph{nērākap pōṅkaḷ. piṟaku iṭatu pakkam tirumpuṅkaḷ.}
  \item \textquotedblleft{}Opposite the temple there is a bank\textquotedblright{} — \emph{kōyilukku etirē oru vaṅki irukkiṟatu}
  \item \textquotedblleft{}Turn right at the second street\textquotedblright{} — \emph{iraṇṭāvatu teruvil valatu pakkam tirumpuṅkaḷ}
  \item the way from your door to the nearest shop, in three requests
\end{itemize}

\begin{tcolorbox}[breakable,colback=teal!4,colframe=teal!35!black,title={Before you move on}]
\textquotedblleft{}Go straight on\textquotedblright{}? (\textbf{\ta{நேராகப்} \ta{போங்கள்}}.) \textquotedblleft{}Turn left\textquotedblright{}? (\textbf{\ta{இடது} \ta{பக்கம்} \ta{திரும்புங்கள்}}.) \textquotedblleft{}Opposite the temple\textquotedblright{}? (\textbf{\ta{கோயிலுக்கு} \ta{எதிரே}}.) Which form asks, and which one waits for another verb? (\textbf{\ta{திரும்புங்கள்}} asks; \textbf{\ta{திரும்பி}} waits.)
\end{tcolorbox}
\section[(review)]{Review — asking the way}

\label{lesson:TA-R288-asking-the-way}



\begin{quote}
Ask how to get to the hospital. Then read \textbf{\ta{ஈரம்}} aloud, and name its first letter. (\textbf{\ta{ஈ}}, \emph{ī}.)
\end{quote}

\subsection*{Your turn}
Say these aloud:

\begin{itemize}
\raggedright
  \item Excuse me, where is the bank? How far is it from here?
  \item Go straight on, then turn right. It is opposite the temple.
  \item It is near. You can walk.
  \item Because it rained, the street is wet. Go slowly.
  \item the way from the bus station to your house, for a visitor
\end{itemize}

\begin{tcolorbox}[breakable,colback=teal!4,colframe=teal!35!black,title={Before you move on}]
\textquotedblleft{}How do I get to …?\textquotedblright{} (\textbf{… -\ta{க்கு} \ta{எப்படிப்} \ta{போக} \ta{வேண்டும்}?}) \textquotedblleft{}How far?\textquotedblright{} (\textbf{\ta{எவ்வளவு} \ta{தூரம்}?}) \textquotedblleft{}Turn left\textquotedblright{}? (\textbf{\ta{இடது} \ta{பக்கம்} \ta{திரும்புங்கள்}}.) \textquotedblleft{}Between the shop and the school\textquotedblright{}? (\textbf{\ta{கடைக்கும்} \ta{பள்ளிக்கும்} \ta{இடையில்}}.)
\end{tcolorbox}
